\Needspace{12\baselineskip}
\originalchapter{相对论电子与自旋回波}
\originalsection{磁谱仪和速度选择}
Junior Lab相对论动力学装置用磁场选择电子动量，再用相互垂直的电磁场选择速度，并由半导体探测器测量沉积能量。源手册采用高斯单位，本册统一用SI：垂直磁场中的曲率关系 $p=eBr$，速度选择条件 $eE=evB$，故 $v=E/B=V/(dB)$。几何狭缝定义的是曲率半径 $r$，平板间距是 $d$，二者不可混用。

相对论关系为 $p=\gamma m_ev$、$K=(\gamma-1)m_ec^2$，等价于 $K=\sqrt{p^2c^2+m_e^2c^4}-m_ec^2$。当 $p\ll m_ec$，展开恢复 $p^2/(2m_e)$；高能时牛顿预测持续偏离。由每一组 $p,v$ 可以算 $m_e=p/(\gamma v)$，但共同磁场校准误差会影响所有点。验证曲线形状和独立估计绝对常数是不同精度目标。

探测器必须标定能量、测量背景并检查电子在窗口的能损。狭缝宽度、磁场非均匀性、平板间距和漏场影响接受速度宽度。将磁场单位gauss直接插入SI公式而未用 $1\,\mathrm G=10^{-4}\,\mathrm T$，会出现四个数量级错误。观察 $K(v)$ 后应比较两种模型的带误差残差，不能只画一条理论线便宣称相对论已被数据核验。

\originalsection{脉冲NMR装置}
公开脉冲NMR材料署名Jordan Kirsch和Robert Newman。静磁场设定Larmor频率 $\omega_0=\gamma B_0$；射频振荡器经开关、功率放大器和线圈产生脉冲，随后样品弱信号经前置放大器、混频器和低通到示波器。混频把射频信号转成与参考频率的差频，便于观察时间包络。接收链在大脉冲后需要恢复，靠近脉冲的盲区不应误当成样品没有磁化。

在旋转坐标系、共振条件下，幅度 $B_1$、时长 $t_p$ 的脉冲使自旋转角 $\theta=\gamma B_1t_p$。先扫描脉宽标定 $\pi/2$ 和 $\pi$ 脉冲。之后测自由感应衰减，既包含真实横向弛豫 $T_2$，又包含静场不均匀导致的相位散开 $T_2^*$。

\originalsection{回波与纵横弛豫}
对静态频率偏移 $\Delta\omega$，$\pi/2$ 后各自相位是 $\Delta\omega t$。在 $t=\tau$ 加一个 $\pi$ 脉冲反转相位方向，第二段积累使 $t=2\tau$ 时重新聚集。这个几何解释消除了静态失谐，不能消除随机不可逆相位涨落，所以回波幅度仍随 $2\tau$ 按 $\e^{-2\tau/T_2}$ 衰减。

纵向弛豫满足 $\dot M_z=(M_0-M_z)/T_1$。理想反转恢复序列给 $M_z(t)=M_0[1-2\e^{-t/T_1}]$，饱和恢复给 $M_0[1-\e^{-t/T_1}]$。重复周期过短会使初态未恢复，拟合模型须相应改变。接收增益、线圈调谐和脉冲误差可同时影响幅度，基线应在拟合中估计，而非任意先扣到零。

\begin{center}\begin{tikzpicture}[>=Stealth]
\draw[->] (0,0)--(8,0) node[right]{$t$};\draw[thick] (.6,0)--(.6,1.1) node[above]{$\pi/2$};\draw[thick] (3.6,0)--(3.6,1.1) node[above]{$\pi$};
\draw[thick,main] (6.1,0) .. controls (6.3,0) and (6.45,1) .. (6.6,1) .. controls (6.75,1) and (6.9,0) .. (7.1,0);
\node[above] at (6.6,1.05){回波};\node[below] at (.6,0){$0$};\node[below] at (3.6,0){$\tau$};\node[below] at (6.6,0){$2\tau$};
\end{tikzpicture}\end{center}

\originalsection{磁谱仪接受范围与参数相关}
垂直磁场下Lorentz力改变动量方向而不做功，$|\dot{\mathbf p}|=evB$，圆轨道方向变化率 $v/r$，所以 $pv/r=evB$，得 $p=eBr$；这里无需预先假设 $p=mv$。速度选择器合力零给 $E=vB$，则同一 $B$ 同时进入 $p$ 和 $v$，二者误差相关。

将原始独立测量写成 $\mathbf z=(B,r,V,d)$，再映射到 $p=eBr,v=V/(dB)$。例如共同磁场相对误差使 $\delta p/p=\delta B/B$、$\delta v/v=-\delta B/B$，产生负相关。按每点的 $p,v$ 独立误差拟合质量会丢失这一项。由能量谱再测 $K$ 可形成第三个独立一致性关系，但探测窗口能损与增益校准必须进入。

无量纲变量 $p/(m_ec)$、$v/c$、$K/(m_ec^2)$ 将关系写成 $p/(m_ec)=\beta_v/\sqrt{1-\beta_v^2}$、$K/(m_ec^2)=1/\sqrt{1-\beta_v^2}-1$。在低速区两模型几乎重合，因此实验必须覆盖有足够理论差异、而速度/动量分辨仍能区分的区间。单纯增加低速重复次数未必最能检验相对论曲率。

\originalsection{Bloch方程、旋转系与脉冲}
宏观磁化在静场中的方程为 $\dot{\mathbf M}=\gamma\mathbf M\times\mathbf B- (M_x\hat x+M_y\hat y)/T_2-(M_z-M_0)\hat z/T_1$。没有弛豫时磁化绕场进动，$\omega_0=|\gamma|B_0$。将旋转射频场移到同频旋转系，快速旋转波近似给有效场 $\mathbf B_{\rm eff}=B_1\hat x+(\Delta\omega/\gamma)\hat z$。

共振时有效场沿横向，转角为 $\gamma B_1t_p$；失谐时转轴倾斜、角速度 $\sqrt{(\gamma B_1)^2+\Delta\omega^2}$，同一个脉宽无法对所有失谐自旋都实现精确 $\pi$。短强脉冲通常激发较宽频带，但会增加接收器恢复盲区和功放非线性。因此脉宽标定须在具体样品和共振条件下进行。

\originalsection{回波如何区分可逆散相和真正弛豫}
不同位置静态场偏移使FID为 $M_+(t)=M_+(0)\e^{-t/T_2}\int p(\Delta\omega)\e^{i\Delta\omega t}\dd\Delta\omega$。Gaussian偏移分布使额外包络为 $\e^{-\sigma_\omega^2t^2/2}$，Lorentzian偏移给指数；$T_2^*$ 只是合并测得的短衰减尺度，未必所有样品都有单指数。

理想 $\pi$ 脉冲把已积累相位变为相反号，$t=2\tau$ 时各静态偏移重新同相。若自旋在梯度中扩散，其两段采样的场不同，回波会因扩散额外衰减；若场随时间随机变，也不能完全复相。所以从回波指数求 $T_2$ 时要检查梯度、扩散和脉冲不完美影响。

为拟合 $T_2$，改变 $\tau$ 并以统一时窗读回波振幅，模型为 $A\e^{-2\tau/T_2}+B$。若先取对数，未知基线和低信噪比会造成偏差，较好在原幅度上用带误差非线性拟合。反转恢复测 $T_1$ 要保留信号符号；取绝对值后零附近呈V形，若仍用原有符号公式拟合会错。不同方法所得的 $T_1,T_2,T_2^*$ 对应不同机制。
